\reading{IM-7}{Portfolio Performance Evaluation}{STRIKE FIRST}{5}

\srcnote{Zvi Bodie, Alex Kane and Alan J. Marcus, \textit{Investments}, 12th
Edition (New York, NY: McGraw-Hill, 2020), Chapter 24. Printed as Chapter 8 in
the GARP source volume.}

\begin{imkeybox}[title={The nine objectives in this reading}]
\begin{enumerate}[label=\textbf{\alph*.}]
\item Differentiate between the time-weighted and dollar-weighted returns of a portfolio and describe their appropriate uses.
\item Describe risk-adjusted performance measures, such as Sharpe's measure, Treynor's measure, Jensen's measure (Jensen's alpha), and the information ratio, and identify the circumstances under which the use of each measure is most relevant.
\item Describe the uses for the Modigliani-squared and Treynor's measure in comparing two portfolios and the graphical representation of these measures.
\item Determine the statistical significance of a performance measure using standard error and the t-statistic.
\item Describe style analysis.
\item Explain the difficulties in measuring the performance of actively managed portfolios.
\item Describe performance manipulation and the problems associated with using conventional performance measures.
\item Describe techniques to measure the market timing ability of fund managers with a regression and with a call option model and calculate a manager's return due to market timing.
\item Describe and apply performance attribution procedures, including the asset allocation decision, sector and security selection decision, and the aggregate contribution.
\end{enumerate}
\end{imkeybox}

\begin{imtrapbox}[breakable=false,title={The reading number, and the order of the letters}]
The source notes head this reading \emph{Reading 90}. Reading 90 is IM-8, Hedge
Fund Investment Strategies, which the notes also cover --- correctly labelled 90
--- fourteen pages later. This reading is 89, and it is \textbf{IM-7}. The notes
also run the objectives out of sequence, placing f, g and h before e and i;
they are restored to spine order below.
\end{imtrapbox}

% ------------------------------------------------------------------
\lo{IM-7 a}{Differentiate between the time-weighted and dollar-weighted returns of a portfolio and describe their appropriate uses.}

\begin{imfmlbox}
\[ r_t = \frac{\text{total proceeds}}{\text{initial investment}}
       = \frac{\text{income} + \text{capital gain}}{\text{initial investment}} \]
\end{imfmlbox}

\begin{itemize}
\item \term{Time-weighted return.} Also known as the \term{geometric average}. It
is \term{not affected by cash flows} during the investment periods.
\item \term{Dollar-weighted return.} Also known as the \term{internal rate of
return (IRR)}: the rate at which the present value of cash inflows equals the
present value of cash outflows.
\end{itemize}

\gapbadge
\srcnote{The notes state the definitions and a one-line rule of thumb but work no
example, so the mechanism behind the difference never appears. GARP's own case
is below (Chapter 8), and every figure in it has been recomputed.}

\begin{imexbox}[title={Why the two averages differ}]
A share costs \$50 and pays a \$2 dividend at year end, when it is worth \$53. A
second share is bought at \$53. At the end of year two both shares pay \$2 and are
sold at \$54 each.

\medskip
\textbf{Time-weighted (geometric):}
\[ r_1 = \frac{53 + 2 - 50}{50} = 10.00\%, \qquad
   r_2 = \frac{54 + 2 - 53}{53} = 5.66\% \]
\[ r_G = (1.10 \times 1.0566)^{1/2} - 1 = \mathbf{7.81\%} \]

\textbf{Dollar-weighted (IRR):} equate the present values of flows,
\[ 50 + \frac{53}{1+r} = \frac{2}{1+r} + \frac{112}{(1+r)^2}
   \;\Longrightarrow\; r = \mathbf{7.117\%} \]

\medskip
The dollar-weighted average is \emph{lower} because the second year's return of
5.66\% --- the worse of the two --- is the year in which \emph{more money was
invested}.
\end{imexbox}

\begin{imkeybox}[title={Which to use, and why}]
Time-weighted returns are unaffected by the size and timing of cash flows, so
they isolate the decisions the \term{manager} actually made. Dollar-weighted
returns weight each period by the money at work in it, so they describe what the
\term{investor} actually earned. GARP ranks funds within a comparison universe
using the (usually time-weighted) average.
\end{imkeybox}

\begin{imtrapbox}[breakable=false,title={The rule of thumb in the notes needs its reason}]
The notes compress this to: \emph{specific time performance of a manager, use
dollar-weighted; overall performance of a manager, use time-weighted.} That rule
is only safe if you carry the reason with it. A manager who does not control the
timing of contributions should not be judged on a measure that is driven by
them --- which is why manager \emph{rankings} are time-weighted. Reach for the
dollar-weighted figure when the question is about the money, not the manager.
\end{imtrapbox}

% ------------------------------------------------------------------
\lo{IM-7 b}{Describe risk-adjusted performance measures, such as Sharpe's measure, Treynor's measure, Jensen's measure (Jensen's alpha), and the information ratio, and identify the circumstances under which the use of each measure is most relevant.}

\renewcommand{\arraystretch}{1.4}
\noindent\begin{tabularx}{\textwidth}{@{}>{\raggedright\arraybackslash}p{2.3cm}cX@{}}
\toprule
\textbf{\sffamily\small Measure} & \textbf{\sffamily\small Formula} & \textbf{\sffamily\small When it is most relevant} \\
\midrule
Sharpe ratio & $S_A = \dfrac{\overline{R_A} - \overline{R_F}}{\sigma_A}$ &
Uses standard deviation --- \term{total risk}. Considers both overall return and
diversification. The measure for a portfolio held \emph{on its own}. \\
Treynor measure & $T_A = \dfrac{\overline{R_A} - \overline{R_F}}{\beta_A}$ &
Uses beta --- \term{systematic risk}. Useful for well-diversified portfolios,
where total and systematic risk are similar. Can give different rankings from
Sharpe for \emph{less} diversified portfolios. \\
Jensen's alpha & $\alpha_A = R_A - E(R_A)$ &
The difference between actual return and the return expected from its systematic
risk under the CAPM. A positive and statistically significant alpha indicates
superior performance. \term{Does not consider diversification.} \\
Information ratio & $\text{IR}_A = \dfrac{\overline{R_A} - \overline{R_B}}{\sigma_{A-B}}$ &
Like Sharpe but against a \term{benchmark} instead of the risk-free rate.
Measures risk-adjusted return relative to a specific benchmark. \\
\bottomrule
\end{tabularx}
\renewcommand{\arraystretch}{1}
\figcap{Four measures, three denominators and two numerators. Sharpe and Treynor
share a numerator and differ in the denominator; Sharpe and the information ratio
share the \emph{shape} and differ in what the return is measured against. Used
together, Sharpe and Treynor provide insights into diversification: when they
rank two portfolios differently, the difference is diversification.}

% ------------------------------------------------------------------
\lo{IM-7 c}{Describe the uses for the Modigliani-squared and Treynor's measure in comparing two portfolios and the graphical representation of these measures.}

\begin{imfmlbox}
\[ M^2 = \underbrace{R_F + \left(\frac{R_P - R_F}{\sigma_P}\right)\sigma_M}_{\textstyle \text{the portfolio levered to the market's risk}} - \; \text{market return} \]
\vk{the bracketed term is the \term{Sharpe ratio} of the portfolio. Multiplying it
by $\sigma_M$ and adding $R_F$ produces the return the portfolio would have earned
if it had been scaled to exactly the market's volatility. $M^2$ is how much better
or worse that is than the market itself. A positive $M^2$ means the portfolio
outperformed the market after adjusting for risk.}
\end{imfmlbox}

\figwrap{%
\begin{tikzpicture}
\begin{axis}[frmaxis,
  width=11.0cm, height=5.8cm,
  xmin=0, xmax=0.30, ymin=0.02, ymax=0.155,
  xlabel={standard deviation $\sigma$},
  ylabel={expected return},
  xtick={0,0.15,0.24}, xticklabels={0,$\sigma_M$,$\sigma_P$},
  ytick={0.03,0.095,0.11,0.125},
  yticklabels={$R_F$,market,$M^2$ adj.,$R_P$},
]
\addplot[FRMNavy, very thick, domain=0:0.30, samples=2] {0.03 + 0.4333*x};
\addplot[FRMRule, thick, dashed, domain=0:0.30, samples=2] {0.03 + 0.4333*x*0};
\addplot[only marks, mark=*, mark size=2.2pt, FRMRed] coordinates {(0.24,0.125)};
\addplot[only marks, mark=square*, mark size=2pt, FRMGreen] coordinates {(0.15,0.095)};
\addplot[only marks, mark=*, mark size=2.2pt, FRMGold] coordinates {(0.15,0.11)};
\draw[FRMGold, line width=1pt, {Stealth[length=4pt]}-{Stealth[length=4pt]}]
  (axis cs:0.153,0.095) -- (axis cs:0.153,0.11);
\draw[black!35, thin, dotted] (axis cs:0.15,0.02) -- (axis cs:0.15,0.11);
\draw[black!35, thin, dotted] (axis cs:0.24,0.02) -- (axis cs:0.24,0.125);
\node[font=\scriptsize\sffamily, fill=white, inner sep=1.5pt, anchor=west,
      text=FRMRed] at (axis cs:0.185,0.137) {portfolio $P$};
\node[font=\scriptsize\sffamily, fill=white, inner sep=1.5pt, anchor=west,
      text=FRMGold] at (axis cs:0.168,0.112) {$M^2$};
\node[font=\scriptsize\sffamily, fill=white, inner sep=1.5pt, anchor=east,
      text=FRMGreen, anchor=north east] at (axis cs:0.146,0.0915) {market};
\node[font=\scriptsize\sffamily, fill=white, inner sep=1.5pt, anchor=west,
      text=FRMNavy, rotate=17] at (axis cs:0.03,0.048) {slope $=$ Sharpe ratio of $P$};
\end{axis}
\end{tikzpicture}}%
{$M^2$ read off a graph. The navy line is the set of portfolios you could build by
mixing $P$ with the risk-free asset; its slope \emph{is} $P$'s Sharpe ratio.
De-lever $P$ down to the market's volatility $\sigma_M$ and read off the return:
that is the bracketed term in the formula. $M^2$ is the gold gap between that
point and the market itself. Because the whole construction is a rescaling along
one line, \term{$M^2$ and the Sharpe ratio always rank portfolios identically} ---
$M^2$ simply restates the Sharpe ratio in percentage-return units, which is what
makes it easier to talk about.}

\begin{imkeybox}[title={Ranking portfolios using these ratios}]
\begin{enumerate}
\item \term{Jensen's alpha and the Treynor measure give the same rankings} ---
both divide by, or are defined against, \emph{beta}.
\item \term{Sharpe and $M^2$ give the same rankings} --- both are built on
\emph{total} risk.
\end{enumerate}
The two pairs need not agree with each other, and where they disagree, the
disagreement is about diversification.
\end{imkeybox}

% ------------------------------------------------------------------
\lo{IM-7 d}{Determine the statistical significance of a performance measure using standard error and the t-statistic.}

A positive alpha indicates superior performance, a negative alpha inferior
performance, and a zero alpha normal performance matching the benchmark. To
assess a manager's ability to generate alpha we conduct a $t$-test:

\begin{imkeybox}
$H_0$: true alpha is zero. \qquad $H_A$: true alpha is not zero.
\end{imkeybox}

\begin{imfmlbox}
\[ t = \frac{\alpha - 0}{\text{standard error}(\alpha)} = \frac{\alpha}{\sigma/\sqrt{N}} \]
\vk{$N$ is the number of observations. The $\sigma$ in the denominator is the
standard deviation of the regression \term{residuals} --- the tracking error ---
not the portfolio's total volatility.}
\end{imfmlbox}

\begin{imtrapbox}[breakable=false,title={Which sigma, and why $\sqrt{N}$}]
Substituting the portfolio's own standard deviation for the residual standard
deviation will produce a $t$-statistic that is too small and a manager who looks
insignificant when she is not. The $\sqrt{N}$ is the reason significance is so
hard to establish in practice: quadrupling the sample only doubles the
$t$-statistic, which is why alpha is, as Ang puts it in IM-3, often hard to detect
statistically.
\end{imtrapbox}

% ------------------------------------------------------------------
\lo{IM-7 e}{Describe style analysis.}

Style analysis regresses fund returns on the returns of a set of \term{tradable}
style indices, with the coefficients constrained to be non-negative and to sum to
one, so that they read as portfolio weights. The fitted combination is the fund's
implied style benchmark, and what is left over is the manager's selection
contribution.

\begin{imkeybox}
This is the same machinery as IM-3 f. There, style analysis solved the problem of
\emph{time-varying factor exposures}; here it solves the problem of \emph{finding
the right benchmark} in the first place. One technique, two objectives, in two
different readings.
\end{imkeybox}

% ------------------------------------------------------------------
\lo{IM-7 f}{Explain the difficulties in measuring the performance of actively managed portfolios.}

Hedge fund performance evaluation is complicated because:

\begin{enumerate}
\item Hedge fund risk is \term{not constant over time} (nonlinear risk).
\item Hedge fund holdings are often \term{illiquid} (data smoothing).
\item Hedge fund sensitivity with traditional markets \term{increases in times of
a market crisis and decreases in times of market strength}.
\end{enumerate}

\begin{imtrapbox}[breakable=false,title={The third one is the dangerous one}]
Correlation rising exactly when it hurts is the opposite of what a diversifier is
bought for. A measure calibrated in calm markets will therefore overstate the
diversification the fund actually delivers. Note the direction: sensitivity
\emph{up} in crisis, \emph{down} in strength.
\end{imtrapbox}

% ------------------------------------------------------------------
\lo{IM-7 g}{Describe performance manipulation and the problems associated with using conventional performance measures.}

\begin{enumerate}[label=\alph*)]
\item The Sharpe ratio assumes \term{consistent risk and return}. This works well
for passive strategies with stable profiles.
\item However, actively managed portfolios can change their risk profile ---
becoming more or less risky --- over time.
\item The Sharpe ratio does not capture these dynamic changes and can be
misleading. It might suggest good performance even if the overall portfolio
underperforms the market due to the risk shift.
\end{enumerate}

\begin{imkeybox}[title={The distinction the notes are careful about}]
The issue is \term{not performance manipulation} but rather the
\term{inapplicability of a static measure} (the Sharpe ratio) to a dynamic
situation (actively managed portfolios). The measure is not being gamed; it is
being applied where its assumptions do not hold.
\end{imkeybox}

% ------------------------------------------------------------------
\lo{IM-7 h}{Describe techniques to measure the market timing ability of fund managers with a regression and with a call option model and calculate a manager's return due to market timing.}

\begin{imdefbox}[title={Market timing}]
The ability to predict the future direction of the market and shift funds between
the market index portfolio and risk-free assets, depending on whether the market
will outperform the risk-free asset.
\end{imdefbox}

\subsection{The regression approach}

Assuming \term{no} market timing and a constant beta, the security characteristic
line is a straight line with a constant slope. The timing model instead assumes
the beta can take only two values, a large one and a smaller one:

\begin{imfmlbox}
\[ R_P - R_F = \alpha + \beta_P (R_M - R_F) + M_P (R_M - R_F)\,D + \varepsilon_P \]
\vk{$D$ is a dummy variable: $D = 1$ when $r_M > r_f$, otherwise $D = 0$. The
portfolio's beta is therefore $\beta_P$ in a bear market and $\beta_P + M_P$ in a
bull market. \term{A positive $M_P$ is an indication of the fund manager's timing
ability.}}
\end{imfmlbox}

\figwrap{%
\begin{tikzpicture}
\begin{axis}[frmaxis, name=l1, scale only axis,
  width=5.0cm, height=4.0cm,
  title={No market timing},
  xlabel={$R_M - R_F$}, ylabel={$R_P - R_F$},
  xmin=-0.12, xmax=0.12, ymin=-0.10, ymax=0.14,
  xtick={0}, xticklabels={0}, ytick={0}, yticklabels={0},
]
\addplot[FRMNavy, very thick, domain=-0.12:0.12, samples=2] {0.8*x};
\draw[black!45, thin] (axis cs:-0.12,0) -- (axis cs:0.12,0);
\node[font=\scriptsize\sffamily, fill=white, inner sep=1.5pt, anchor=north west,
      text=FRMNavy] at (axis cs:-0.115,0.135) {constant slope $\beta_P$};
\end{axis}
\begin{axis}[frmaxis, name=l2, at={(l1.right of south east)}, xshift=1.7cm,
  scale only axis, width=5.0cm, height=4.0cm,
  title={With market timing},
  xlabel={$R_M - R_F$},
  xmin=-0.12, xmax=0.12, ymin=-0.10, ymax=0.14,
  xtick={0}, xticklabels={0}, ytick={0}, yticklabels={0},
]
\addplot[FRMNavy, very thick, domain=-0.12:0, samples=2] {0.5*x};
\addplot[FRMNavy, very thick, domain=0:0.12, samples=2] {1.2*x};
\draw[black!45, thin] (axis cs:-0.12,0) -- (axis cs:0.12,0);
\node[font=\scriptsize\sffamily, fill=white, inner sep=1.5pt, anchor=north west,
      text=FRMGreen] at (axis cs:-0.115,0.135) {$\beta_P + M_P$ when $R_M > R_F$};
\node[font=\scriptsize\sffamily, fill=white, inner sep=1.5pt, anchor=south west,
      text=FRMRed] at (axis cs:-0.115,-0.095) {$\beta_P$ when $R_M < R_F$};
\end{axis}
\end{tikzpicture}}%
{The kink is the timing. With no timing ability the security characteristic line
is straight: one beta in all states. A successful timer holds more market
exposure when the market beats the risk-free rate and less when it does not,
which bends the line upward at the origin. The upward bend is exactly the payoff
shape of a \emph{long call}, and estimating $M_P$ is estimating the size of that
bend. Note that this is the same L-shape argument as IM-3 g, running the other
way: there a kink created a spurious alpha, here a kink is the genuine skill being
measured.}

\subsection{The call option model}

The key to valuing market timing ability is to recognise that \term{perfect
foresight is equivalent to holding a call option on the equity portfolio}. The
perfect timer invests 100\% in either the safe asset or the equity portfolio,
whichever will provide the higher return. The value of that perfect foresight,
and therefore the return due to market timing, is the value of the call.

% ------------------------------------------------------------------
\lo{IM-7 i}{Describe and apply performance attribution procedures, including the asset allocation decision, sector and security selection decision, and the aggregate contribution.}

Performance attribution identifies the sources of value added to the portfolio.

\begin{itemize}
\item \term{Asset allocation:} how much of the performance is attributable to the
selection of the risky asset classes.
\item \term{Security selection:} how much is attributable to selection of the
right sector or security within an asset class.
\end{itemize}

\begin{imfmlbox}
\begin{align*}
\text{Contribution from asset allocation} &= (w_P - w_B) \times R_B \\
\text{Contribution from security selection} &= w_P \times (R_P - R_B) \\[3pt]
\text{Total contribution} &= w_P R_P - w_B R_B
\end{align*}
\vk{$w_P$ and $w_B$ are the portfolio's and benchmark's weights in the asset
class; $R_P$ and $R_B$ are the returns earned in it by the portfolio and the
benchmark.}
\end{imfmlbox}

\begin{imkeybox}[title={The identity, and why it is worth checking}]
The two contributions sum exactly to the total:
\[ (w_P - w_B)R_B + w_P(R_P - R_B)
   = w_P R_B - w_B R_B + w_P R_P - w_P R_B
   = w_P R_P - w_B R_B \]
The $w_P R_B$ terms cancel. If a candidate answer's two components do not add to
$w_P R_P - w_B R_B$, one of them is wrong, and you can find which without
recomputing either.
\end{imkeybox}

\begin{imtrapbox}[title={Consolidated trap summary --- IM-7}]
\begin{itemize}
\item \term{Sibling, IM-7 a.} Time-weighted is the geometric average and ignores
cash flows; dollar-weighted is the IRR and is driven by them. The dollar-weighted
figure is lower here precisely because the worse year had more money in it.
\item \term{Sibling, IM-7 b.} Sharpe divides by $\sigma$, Treynor by $\beta$,
the information ratio by tracking error. Jensen's alpha is a difference, not a
ratio, and does not consider diversification.
\item \term{Ranking, IM-7 c.} Jensen and Treynor rank alike; Sharpe and $M^2$ rank
alike. Pairing Sharpe with Treynor, or Jensen with $M^2$, is the corruption.
\item \term{Formula, IM-7 c.} $M^2$ subtracts the \emph{market} return, not the
risk-free rate. Forgetting the subtraction leaves the risk-adjusted portfolio
return, which is a real and offerable number.
\item \term{Definition, IM-7 d.} The $\sigma$ in the $t$-statistic is the residual
standard deviation, not the portfolio's volatility.
\item \term{Scope, IM-7 g.} The Sharpe ratio problem with active portfolios is
\emph{inapplicability of a static measure}, not manipulation.
\item \term{Polarity, IM-7 f.} Hedge fund sensitivity to traditional markets rises
in a crisis and falls in strength --- the unhelpful direction.
\item \term{Sign, IM-7 h.} A \emph{positive} $M_P$ indicates timing ability, and
the dummy is $D = 1$ when $r_M > r_f$. Reversing the dummy reverses the
conclusion.
\item \term{Intermediate result, IM-7 i.} Allocation uses the \emph{benchmark's}
return $R_B$; selection uses the \emph{portfolio's} weight $w_P$. Swapping which
is which still produces two numbers that look right individually and fail the
sum check.
\end{itemize}
\end{imtrapbox}

% ==================================================================
